
\chapter{Cadre général et fondements du projet}

\section*{Introduction}
\addcontentsline{toc}{section}{Introduction}

Dans ce chapitre, nous présentons le cadre général du projet, son contexte professionnel, organisationnel et technologique. Nous présentons également l'entreprise d'accueil, le département dans lequel le stage a été réalisé, le contexte ayant motivé le sujet, puis les fondements nécessaires à la compréhension de la solution proposée.

\section{Présentation de l'entreprise d'accueil}

Dans cette section, nous présentons l'entreprise d'accueil Talan Tunisie Consulting, ses principaux secteurs d'activité, le département Tech et Cloud ainsi que le bootcamp d'intégration suivi au début du stage. Cette présentation permet de situer le projet dans son environnement professionnel et de mieux comprendre le cadre dans lequel il a été réalisé.

\subsection{Talan Tunisie Consulting}

Talan est un groupe international de conseil et d'expertise technologique fondé en 2002. Il accompagne les entreprises et les institutions publiques dans leurs projets de transformation digitale, d'innovation et d'amélioration de leurs processus métier \cite{talan_org}. Son activité se situe à l'intersection du conseil, de l'ingénierie logicielle, de la data, du cloud et de l'accompagnement des organisations dans l'évolution de leurs systèmes d'information.

Le groupe est présent sur cinq continents et dans 21 pays. Il s'appuie sur plusieurs centres d'excellence, notamment à Tunis, Luxembourg, Île Maurice, Hongrie, Espagne, Pologne et Mexique. Cette organisation internationale lui permet de combiner proximité client, expertise technologique et capacité de réalisation à distance, selon les besoins de chaque mission \cite{talan_org}.

Talan compte plus de 7\,000 collaborateurs. En 2024, le groupe a réalisé un chiffre d'affaires de 850 millions d'euros et vise à dépasser le milliard d'euros à l'horizon 2026 \cite{talan_org}. Au-delà de sa croissance, l'entreprise met en avant une culture orientée vers l'humain, l'apprentissage continu et l'innovation, comme l'illustrent ses reconnaissances Great Place To Work \cite{gptw_talan}.

Talan Tunisie s'inscrit dans cette dynamique globale. Elle constitue un pôle d'expertise important pour la réalisation de projets technologiques, l'accompagnement des clients internationaux et le développement de solutions logicielles répondant à des problématiques métier variées. Le contexte de travail offert par cette structure a été particulièrement adapté à un projet de fin d'études portant sur l'architecture logicielle, l'automatisation et la modernisation d'applications existantes.

\subsection{Secteurs d'activité}

Talan intervient dans plusieurs secteurs d'activité, notamment l'énergie, le transport et la logistique, la banque et la finance, l'industrie, l'assurance et le service public \cite{talan_secteurs}. Ces secteurs ont des besoins différents, mais ils partagent souvent des enjeux communs : moderniser des applications existantes, améliorer la maintenabilité, accélérer la livraison de nouvelles fonctionnalités et réduire les risques liés aux évolutions techniques.

Afin de mieux présenter la diversité des domaines d'intervention de Talan, le tableau~\ref{tab:secteurs-talan} regroupe quelques secteurs représentatifs ainsi que les enjeux généralement associés à chacun d'eux.

\begingroup
\renewcommand{\arraystretch}{1.25}
\begin{longtable}{|p{3.6cm}|p{10.4cm}|}
\caption{Exemples de secteurs couverts par Talan et enjeux associés}\label{tab:secteurs-talan}\\
\hline
\textbf{Secteur} & \textbf{Enjeux généralement adressés} \\
\hline
Énergie & Transformation des systèmes, suivi de la performance, intégration de solutions numériques et optimisation des processus opérationnels. \\
\hline
Transport et logistique & Amélioration de la planification, fiabilisation des flux, modernisation des plateformes et meilleure exploitation des données. \\
\hline
Banque et finance & Évolution des services digitaux, sécurisation des applications, automatisation des processus et amélioration de l'expérience utilisateur. \\
\hline
Industrie & Digitalisation des chaînes de valeur, suivi des indicateurs, intégration de systèmes et amélioration de la qualité logicielle. \\
\hline
Assurance & Modernisation des parcours clients, simplification des outils internes et meilleure interconnexion des services. \\
\hline
Service public & Transformation numérique, rationalisation des démarches et amélioration de l'accessibilité des services. \\
\hline
\end{longtable}
\endgroup

Cette diversité sectorielle montre que les projets confiés à Talan ne se limitent pas à la mise en place d'outils techniques. Ils impliquent aussi une compréhension du contexte métier, des contraintes de maintenance, des exigences de sécurité et des objectifs de transformation propres à chaque organisation. Dans ce type d'environnement, les architectures logicielles jouent un rôle essentiel, car elles conditionnent la capacité d'un système à évoluer durablement.

\subsection{Département Tech et Cloud}

Le projet a été réalisé au sein du département \textbf{Tech et Cloud} de Talan Tunisie. Ce département intervient sur des missions liées au développement logiciel, à l'architecture applicative, au cloud computing, à l'industrialisation des processus de développement et à la modernisation des systèmes existants. Il accompagne les clients dans la conception de solutions plus fiables, maintenables et adaptées aux contraintes de production.

Ses activités couvrent plusieurs dimensions complémentaires. La première concerne la conception et le développement d'applications, avec une attention particulière portée à la qualité du code, à la structuration des projets et à la cohérence des choix techniques. La deuxième dimension concerne l'architecture, c'est-à-dire la manière d'organiser les composants logiciels, les échanges, les dépendances et les mécanismes de déploiement. La troisième dimension est liée au cloud et au DevOps, avec des pratiques orientées automatisation, intégration continue, déploiement continu et observabilité.

Ce cadre a été pertinent pour notre sujet, car la migration d'une application front-end monolithique vers une architecture plus modulaire ne relève pas uniquement du développement. Elle demande également une réflexion d'architecture, une capacité d'analyse du code existant, une maîtrise de l'automatisation et une attention particulière à la validation du résultat. Le département Tech et Cloud a donc fourni un environnement cohérent avec les exigences du projet.

\subsection{Bootcamp d'intégration}

Au début du stage, nous avons participé à un bootcamp d'intégration de deux semaines. Cette phase avait pour objectif de faciliter l'entrée dans l'environnement de Talan, de renforcer les compétences nécessaires au projet et d'introduire plusieurs notions liées à l'utilisation de l'intelligence artificielle dans le cycle de vie logiciel.

La première semaine était consacrée à l'intégration de l'intelligence artificielle dans les différentes phases du développement logiciel. Les travaux abordés concernaient notamment la génération de code, l'aide à la documentation, l'automatisation des tests, l'assistance à la revue de code et l'intégration de ces pratiques dans des pipelines de développement. Cette partie a permis de comprendre l'intérêt de l'IA comme levier d'accélération, mais aussi les précautions nécessaires en matière de qualité, de sécurité et de contrôle.

La deuxième semaine portait sur les systèmes Agentic AI et multi-agents. L'accent a été mis sur l'orchestration des tâches, la spécialisation des agents, la validation des sorties, la gestion du contexte, l'observabilité et la maîtrise des droits d'accès. Cette partie a été directement utile pour notre projet, car la solution proposée repose sur une orchestration de plusieurs traitements spécialisés, coordonnés autour d'un objectif commun : produire une migration fiable et exploitable.

Pour clarifier l'apport de cette phase d'intégration, le tableau~\ref{tab:bootcamp-apports} présente une synthèse des deux semaines du bootcamp et de leur contribution directe au projet.

\begingroup
\renewcommand{\arraystretch}{1.25}
\begin{longtable}{|p{3.2cm}|p{5.5cm}|p{5.5cm}|}
\caption{Apports du bootcamp par rapport au projet}\label{tab:bootcamp-apports}\\
\hline
\textbf{Période} & \textbf{Contenu principal} & \textbf{Apport pour le projet} \\
\hline
Semaine 1 & IA appliquée au cycle de vie du développement logiciel & Compréhension de l'utilisation de l'IA pour assister l'analyse, la génération, les tests et la documentation. \\
\hline
Semaine 2 & Systèmes Agentic AI et orchestration multi-agents & Préparation à la conception d'un pipeline composé d'étapes spécialisées, contrôlées et validées progressivement. \\
\hline
\end{longtable}
\endgroup

Ce bootcamp n'a pas remplacé le travail de conception du projet, mais il a fourni une base méthodologique et technique. Il a permis de mieux identifier les points sensibles d'un système automatisé : la définition des responsabilités, la transmission des informations entre étapes, la limitation des erreurs et la nécessité de produire des résultats vérifiables.

\section{Présentation du projet}

Dans cette section, nous présentons le contexte du projet, le constat initial, la problématique, les objectifs ainsi que le périmètre retenu. Cette présentation permet de clarifier le besoin traité, de situer les difficultés rencontrées et de préciser les limites de la solution proposée.

\subsection{Contexte du projet}

Le projet s'inscrit dans une démarche de modernisation d'une application front-end Angular initialement organisée sous forme monolithique. Dans ce type de contexte, l'application peut devenir progressivement difficile à faire évoluer lorsque le nombre de fonctionnalités augmente, lorsque plusieurs intervenants doivent travailler sur le même code ou lorsque les dépendances internes deviennent nombreuses.

L'objectif du stage n'était pas de réaliser une simple application Angular ni de concevoir manuellement une architecture cible. Le besoin portait sur l'automatisation du processus de transformation : partir d'un projet existant, analyser sa structure, proposer un découpage cohérent, générer une nouvelle organisation et produire un résultat utilisable par un développeur. Les définitions détaillées des architectures concernées sont donc volontairement placées dans l'état de l'art, tandis que cette section précise uniquement le contexte opérationnel du projet.

Le point de départ est un projet Angular monolithique exécutable. Il contient des composants, des routes, des services, des dépendances, des styles et des éléments partagés. Le point d'arrivée attendu est un projet organisé selon une architecture micro-frontend, avec une application principale et plusieurs sous-applications. Entre ces deux états, plusieurs décisions doivent être prises : quelles parties regrouper, quelles dépendances partager, quelles routes déplacer et comment vérifier que le projet généré reste cohérent.

\subsection{Constat initial}

Avant de définir la solution, un constat a été établi à partir des difficultés généralement rencontrées lors de la modernisation d'applications front-end existantes. Une migration manuelle demande beaucoup de temps, car elle oblige le développeur à parcourir la structure du projet, comprendre les liens entre fichiers, identifier les dépendances, déplacer les éléments et adapter les configurations. Chaque étape peut introduire des erreurs difficiles à détecter tardivement.

Afin de structurer ce constat, le tableau~\ref{tab:constat-projet} présente les principaux problèmes identifiés ainsi que leurs conséquences sur le processus de migration.

\begingroup
\renewcommand{\arraystretch}{1.25}
\begin{longtable}{|p{4cm}|p{10cm}|}
\caption{Constats ayant motivé le projet}\label{tab:constat-projet}\\
\hline
\textbf{Constat} & \textbf{Conséquence sur la migration} \\
\hline
Taille croissante du code front-end & L'analyse manuelle devient plus longue et plus exposée aux oublis. \\
\hline
Dépendances internes nombreuses & Le déplacement d'un composant ou d'un service peut provoquer des erreurs indirectes. \\
\hline
Routage centralisé & La séparation des parcours fonctionnels demande une adaptation rigoureuse. \\
\hline
Manque de traçabilité & Il devient difficile de justifier les choix de migration sans documentation générée. \\
\hline
Validation tardive & Les erreurs détectées uniquement à la fin du processus coûtent plus cher à corriger. \\
\hline
\end{longtable}
\endgroup

Ces constats montrent que le problème ne se limite pas au déplacement de fichiers. Il concerne l'ensemble de la chaîne de transformation, depuis la compréhension de l'application source jusqu'à la vérification du projet produit.

\subsection{Problématique}

La transformation d'une application Angular monolithique en architecture micro-frontend ne consiste pas seulement à séparer le code en plusieurs projets. Elle nécessite d'analyser l'existant, d'identifier des domaines fonctionnels cohérents, de définir une architecture cible, puis de migrer correctement les composants, les services, les routes, les dépendances et les mécanismes de communication.

La difficulté principale réside dans la maîtrise de cette chaîne de transformation. Un découpage mal défini peut produire une architecture trop fragmentée, difficile à maintenir, ou au contraire une architecture trop proche du monolithe initial. Une configuration incomplète peut aussi empêcher le build ou l'exécution du projet généré. Enfin, l'absence de documentation peut rendre le résultat difficile à comprendre et à reprendre par un développeur.

La problématique de ce projet peut donc être formulée ainsi : \textit{comment automatiser la migration d'un projet front-end Angular monolithique vers une architecture micro-frontend exécutable, tout en contrôlant la qualité des transformations et en limitant les erreurs de migration ?}

Cette problématique met en évidence trois dimensions. La première est l'analyse du projet source. La deuxième est la génération d'une architecture cible cohérente. La troisième est la validation progressive du résultat, afin que la migration produise un livrable exploitable et documenté.

\subsection{Objectifs du projet}

Les objectifs du projet ont été définis de manière à couvrir l'ensemble du processus de migration. Ils ne concernent pas uniquement la génération de fichiers, mais aussi la compréhension de l'application source, la structuration de l'architecture cible, la validation des résultats et la production d'une documentation utile.

Pour rendre ces objectifs plus lisibles, le tableau~\ref{tab:objectifs-projet} les regroupe selon les principales actions attendues de la solution.

\begingroup
\renewcommand{\arraystretch}{1.25}
\begin{longtable}{|p{3.3cm}|p{10.7cm}|}
\caption{Objectifs du projet}\label{tab:objectifs-projet}\\
\hline
\textbf{Objectif} & \textbf{Description} \\
\hline
Analyser l'existant & Extraire les informations utiles du projet Angular source : structure, routes, composants, services, dépendances et éléments partagés. \\
\hline
Proposer un découpage & Identifier des domaines fonctionnels pouvant devenir des micro-frontends cohérents. \\
\hline
Générer la cible & Créer une structure de projet organisée autour d'une application principale et de plusieurs applications fonctionnelles. \\
\hline
Migrer les artefacts & Déplacer ou reconstruire les éléments nécessaires : composants, templates, styles, services, routes et configurations. \\
\hline
Valider le résultat & Vérifier la cohérence de l'architecture produite, la présence des fichiers attendus et la capacité du projet à être construit. \\
\hline
Documenter la migration & Produire une documentation indiquant les choix effectués, les transformations appliquées et la structure finale. \\
\hline
\end{longtable}
\endgroup

Ces objectifs servent de base à la planification du projet et au découpage des tâches. Ils permettent également d'évaluer la solution, car chaque objectif correspond à un livrable ou à un résultat observable.

\subsection{Périmètre du projet}

Le périmètre du projet a été défini pour garder un équilibre entre ambition et faisabilité. La solution vise les projets Angular monolithiques et se concentre sur les éléments front-end nécessaires à la migration. Elle prend en compte la structure applicative, les routes, les composants, les services et les dépendances utiles à l'architecture cible.

En revanche, le projet ne vise pas à réécrire automatiquement toute la logique métier d'une application complexe ni à garantir la migration parfaite de tous les cas particuliers possibles. Certains ajustements peuvent rester nécessaires lorsque le projet source contient des spécificités fortes, des bibliothèques peu standardisées ou des intégrations non documentées.

Le tableau~\ref{tab:perimetre-projet} précise les éléments inclus dans le périmètre du projet et ceux qui ne constituent pas l'objectif principal de cette solution.

\begingroup
\renewcommand{\arraystretch}{1.25}
\begin{longtable}{|p{4cm}|p{5cm}|p{5cm}|}
\caption{Délimitation du périmètre du projet}\label{tab:perimetre-projet}\\
\hline
\textbf{Aspect} & \textbf{Inclus dans le périmètre} & \textbf{Hors périmètre principal} \\
\hline
Technologie front-end & Projets Angular et structure associée & Migration vers d'autres frameworks front-end \\
\hline
Architecture cible & Shell, Remotes et configuration de fédération & Refonte complète de l'ergonomie ou du design produit \\
\hline
Migration technique & Composants, routes, services, styles et dépendances & Réécriture métier profonde non déductible du code source \\
\hline
Validation & Contrôles structurels, configuration et build & Audit de sécurité complet ou tests de charge avancés \\
\hline
Documentation & Rapport de migration et structure générée & Documentation fonctionnelle exhaustive de l'application métier \\
\hline
\end{longtable}
\endgroup

Cette délimitation permet de clarifier les attentes : la solution doit assister fortement la migration, produire une base exécutable et réduire l'effort manuel, mais elle ne remplace pas toute expertise d'architecture dans les cas complexes.

\section{Intérêt et motivations du sujet}

Dans cette section, nous présentons l'intérêt industriel, technique et académique du sujet. Nous mettons en évidence les raisons qui justifient le choix de cette problématique, notamment dans un contexte de modernisation applicative et d'automatisation des processus de développement.

Ce sujet présente un intérêt industriel direct. De nombreuses entreprises possèdent des applications front-end développées sur plusieurs années, parfois devenues difficiles à faire évoluer. Une approche automatisée peut réduire l'effort manuel, améliorer la traçabilité et accélérer les premières étapes d'une migration architecturale.

Pour une entreprise de conseil comme Talan, la modernisation applicative représente un enjeu concret. Les clients souhaitent souvent préserver la valeur de leurs systèmes existants tout en les rendant plus évolutifs. La capacité à proposer des démarches structurées, reproductibles et outillées constitue donc un avantage dans les missions d'accompagnement.

Le projet présente également un intérêt technique. Il mobilise des compétences en analyse de code, architecture front-end, génération de projets, orchestration de traitements, validation automatique et documentation. Ces dimensions correspondent à des problématiques actuelles du développement logiciel, où l'automatisation et l'intelligence artificielle sont de plus en plus utilisées pour assister des tâches complexes.

Enfin, le sujet est motivant sur le plan académique, car il oblige à relier plusieurs domaines : architecture logicielle, ingénierie front-end, méthodologie agile et systèmes multi-agents. Il ne s'agit pas seulement de développer une fonctionnalité isolée, mais de concevoir un processus complet, contrôlé et orienté résultat.

\section{Fondements technologiques et état de l'art}

Cette section regroupe les notions techniques nécessaires à la compréhension du projet. Elle présente les architectures front-end, les micro-frontends, la migration d'applications Angular, Module Federation, les grands modèles de langage, le concept d'Agentic AI, l'automatisation et l'orchestration multi-agents.

\subsection{Évolution des architectures front-end}

Les applications web sont passées d'interfaces simples, souvent générées côté serveur, à des applications riches construites avec des frameworks front-end. Cette évolution a amélioré l'interactivité et la réutilisation des composants, mais elle a aussi augmenté la taille du code, le nombre de dépendances et la complexité de maintenance.

Dans ce contexte, l'architecture devient un facteur important de maîtrise du projet. Une application peut rester monolithique lorsque son périmètre est limité. Elle peut aussi être structurée en modules internes, puis évoluer vers des micro-frontends lorsque les domaines fonctionnels, les équipes et les cycles de livraison doivent être mieux séparés.

\subsection{Architecture front-end monolithique}

Une architecture front-end monolithique regroupe l'ensemble de l'interface utilisateur dans une seule application. Les composants, les services, le routage, les styles, la configuration et les dépendances sont maintenus dans un même projet et livrés comme une unité unique.

Cette organisation est souvent pertinente au démarrage d'un projet, car elle simplifie la configuration, le routage et le déploiement. Ses limites apparaissent lorsque le périmètre fonctionnel s'étend : les relations entre fichiers deviennent plus difficiles à suivre, les dépendances internes se multiplient et plusieurs équipes doivent intervenir dans le même cycle de livraison.

Le tableau~\ref{tab:forces-limites-monolithe} résume les principaux atouts d'une architecture front-end monolithique ainsi que les limites qui peuvent motiver une démarche de migration.

\begingroup
\renewcommand{\arraystretch}{1.25}
\begin{longtable}{|p{6.8cm}|p{6.8cm}|}
\caption{Forces et limites d'une architecture front-end monolithique}\label{tab:forces-limites-monolithe}\\
\hline
\textbf{Forces} & \textbf{Limites} \\
\hline
Mise en place rapide et structure initiale simple & Croissance du code plus difficile à maîtriser \\
\hline
Déploiement unique et configuration centralisée & Dépendances internes nombreuses et parfois implicites \\
\hline
Compréhension globale plus facile au démarrage & Risque d'effets de bord lors des modifications \\
\hline
Adapté aux petites équipes et aux périmètres stables & Collaboration plus complexe lorsque plusieurs équipes interviennent \\
\hline
\end{longtable}
\endgroup

Ces limites ne rendent pas le monolithe systématiquement inadapté. Elles indiquent surtout qu'un autre style d'architecture peut devenir nécessaire lorsque l'application doit évoluer plus vite, avec davantage d'autonomie entre domaines fonctionnels.

\subsection{Architecture micro-frontend}

\subsubsection{Définition et principe général}

L'architecture micro-frontend applique au front-end une logique de découpage inspirée des microservices. Elle consiste à diviser une application en plusieurs sous-applications autonomes, chacune associée à un domaine fonctionnel précis et pouvant être développée, testée et maintenue de manière plus indépendante \cite{single_spa_microfrontend}.

Pour l'utilisateur final, l'application reste perçue comme une interface unique. Sur le plan technique, cette interface est composée de plusieurs parties intégrées dans une même expérience globale. La figure~\ref{fig:mfe_independent_deploy} illustre ce principe : plusieurs micro-frontends sont construits séparément, puis assemblés dans une application finale.

\begin{figure}[H]
    \centering
    \includegraphics[width=0.95\textwidth]{figures/mfe.png}
    \caption{Illustration du déploiement indépendant de plusieurs micro-frontends \cite{aymax_microfrontends}}
    \label{fig:mfe_independent_deploy}
\end{figure}

\subsubsection{Fondements d'une architecture micro-frontend : Shell et Remotes}

Dans une architecture micro-frontend, nous distinguons généralement une application hôte, appelée \textit{Shell}, et des applications distantes, appelées \textit{Remotes}. Le \textit{Shell} constitue le point d'entrée global : il porte la navigation principale, charge les micro-frontends et compose l'interface finale présentée à l'utilisateur \cite{single_spa_microfrontend,nx_microfrontend}.

Les \textit{Remotes} représentent les parties fonctionnelles intégrées au \textit{Shell}. Chaque \textit{Remote} doit correspondre à un domaine métier cohérent, par exemple un espace utilisateur, un catalogue ou un module d'administration. La qualité de cette séparation dépend donc des frontières fonctionnelles, mais aussi du routage, du partage des dépendances et de la stratégie de build.

\subsubsection{Module Federation}

Module Federation est une fonctionnalité de Webpack 5 permettant à plusieurs applications construites séparément d'exposer et de consommer des modules au moment de l'exécution \cite{module_federation,webpack_module_federation_docs}. Dans une architecture micro-frontend, elle permet au \textit{Shell} de charger dynamiquement des modules exposés par les \textit{Remotes} sans les intégrer dans un bundle unique au moment de la compilation.

Le principe repose sur deux configurations complémentaires. Côté \textit{Remote}, l'option \texttt{exposes} déclare les modules rendus accessibles via un point d'entrée distant, souvent nommé \texttt{remoteEntry.js}. Côté \textit{Shell}, l'option \texttt{remotes} déclare les applications distantes qui peuvent être chargées \cite{webpack_module_federation_plugin}. Le partage des bibliothèques communes est configuré à travers \texttt{shared}, afin de réduire les duplications et les conflits de versions \cite{module_federation_shared,nx_mf_versions}.

Dans le cas d'Angular, ce partage doit être contrôlé avec attention. Des bibliothèques comme \texttt{@angular/core}, \texttt{@angular/common}, \texttt{@angular/router} et \texttt{rxjs} doivent rester cohérentes entre le \textit{Shell} et les \textit{Remotes}. Le tableau~\ref{tab:module-federation-elements} regroupe les éléments de configuration les plus sensibles pour une application Angular basée sur Module Federation.

\begingroup
\renewcommand{\arraystretch}{1.08}
\begin{longtable}{|p{3.8cm}|p{5.4cm}|p{5.2cm}|}
\caption{Éléments critiques dans une configuration Module Federation}\label{tab:module-federation-elements}\\
\hline
\textbf{Élément} & \textbf{Rôle} & \textbf{Point de validation} \\
\hline
\texttt{remoteEntry.js} & Point d'entrée généré par la \textit{Remote}. & Vérifier l'adresse distante et la disponibilité du fichier. \\
\hline
\texttt{exposes} & Modules rendus accessibles par une \textit{Remote}. & Vérifier que chaque module exposé existe réellement. \\
\hline
\texttt{remotes} & Applications distantes déclarées côté \textit{Shell}. & Vérifier la cohérence entre les noms déclarés côté \textit{Shell} et côté \textit{Remote}. \\
\hline
\texttt{shared} & Dépendances mutualisées entre applications fédérées. & Vérifier le partage des bibliothèques Angular et RxJS. \\
\hline
\texttt{singleton} & Usage d'une seule instance d'une dépendance partagée. & L'activer pour les dépendances Angular critiques. \\
\hline
\texttt{requiredVersion} et \texttt{strictVersion} & Contrôle de compatibilité des versions. & Détecter les versions incompatibles avant l'exécution. \\
\hline
Routes associées & Liaison entre la navigation du \textit{Shell} et les modules distants. & Vérifier que chaque route pointe vers la bonne \textit{Remote}. \\
\hline
\end{longtable}
\endgroup

\subsubsection{Modes d'intégration}

Deux modes d'intégration sont généralement distingués. Dans l'intégration à la compilation, les parties à intégrer sont connues lors de la construction de l'application. Cette approche facilite certains contrôles, mais réduit l'indépendance des modules, car une modification importante peut imposer une reconstruction globale.

Dans l'intégration à l'exécution, le \textit{Shell} charge les \textit{Remotes} dynamiquement, par exemple avec Module Federation. Cette stratégie favorise l'indépendance des livraisons, mais elle impose une configuration plus rigoureuse des points d'exposition, des versions et des dépendances partagées.

Le choix de la granularité reste également important. Une granularité trop fine multiplie les points d'intégration ; une granularité trop large conserve une partie des limites du monolithe. Dans une migration, la granularité doit donc être guidée par les domaines fonctionnels plutôt que par la seule structure des dossiers.

\subsubsection{Points de vigilance}

Les micro-frontends apportent de la modularité, mais ils introduisent aussi des contraintes qu'il faut maîtriser. Les principaux points de vigilance sont synthétisés dans le tableau~\ref{tab:vigilance-mfe}.

\begingroup
\renewcommand{\arraystretch}{1.25}
\begin{longtable}{|p{4cm}|p{10cm}|}
\caption{Points de vigilance dans une architecture micro-frontend}\label{tab:vigilance-mfe}\\
\hline
\textbf{Point de vigilance} & \textbf{Explication} \\
\hline
Frontières fonctionnelles & Les \textit{Remotes} doivent correspondre à des domaines cohérents et non à un découpage arbitraire. \\
\hline
Dépendances partagées & Les bibliothèques communes doivent être gérées pour éviter les conflits et les duplications inutiles. \\
\hline
Routage & La navigation globale et les routes internes doivent rester compréhensibles. \\
\hline
Expérience utilisateur & Les interfaces produites par plusieurs \textit{Remotes} doivent conserver une homogénéité visuelle et comportementale. \\
\hline
Observabilité & Les erreurs doivent pouvoir être localisées dans le \textit{Shell} ou dans la \textit{Remote} concernée. \\
\hline
\end{longtable}
\endgroup

\subsection{Migration vers les micro-frontends}

La migration vers les micro-frontends transforme progressivement une application existante en un ensemble de sous-applications plus autonomes. Elle s'appuie sur deux décisions liées : choisir l'ordre d'extraction des fonctionnalités et définir les frontières fonctionnelles de chaque futur micro-frontend.

\subsubsection{Migration progressive et découpage fonctionnel}

Une migration progressive limite les risques en évitant une réécriture complète. Elle commence souvent par un domaine isolable, dont les routes, composants et services sont clairement identifiables. Par exemple, dans une application de gestion, un espace \og catalogue \fg{} peut être extrait en premier s'il possède ses propres pages, ses services d'accès aux données et peu de dépendances avec les autres domaines.

Le découpage fonctionnel consiste donc à regrouper les éléments du code selon leur rôle métier. Dans une application Angular, cette analyse doit tenir compte des modules, composants, services injectables, routes, dépendances partagées et relations entre les différentes parties de l'application. Un bon découpage évite deux dérives opposées : une fragmentation trop fine, qui alourdit l'intégration, et une séparation trop large, qui reproduit le monolithe initial.

\subsubsection{Risques liés à la migration}

La migration vers une architecture micro-frontend expose le projet à plusieurs niveaux de risques. Sur le plan fonctionnel, une fonctionnalité extraite sans toutes ses dépendances métier peut devenir incomplète ou incohérente. Sur le plan technique, des erreurs peuvent apparaître dans le routage, les imports, les dépendances partagées ou la configuration de build.

Sur le plan architectural, un mauvais découpage peut produire des \textit{Remotes} trop dépendantes les unes des autres, ou au contraire trop nombreuses. Sur le plan opérationnel, le déploiement séparé du \textit{Shell} et des \textit{Remotes} nécessite une compatibilité maîtrisée entre versions. Enfin, la dimension organisationnelle doit être anticipée : lorsque plusieurs équipes interviennent, il faut définir des conventions communes, des responsabilités claires et des règles de cohérence visuelle.

\subsection{Angular : composants, services et routage dans le contexte d'une migration}

Le choix d'Angular dans ce projet n'est pas théorique : la solution vise précisément des applications Angular monolithiques à migrer vers une architecture micro-frontend. Angular est structuré autour de composants, de services, de routes et d'un système d'injection de dépendances \cite{angular_docs}. Ces éléments constituent donc les unités principales à analyser avant toute migration.

Dans une application Angular, un composant ne se résume pas à un fichier TypeScript. Il est généralement lié à un template HTML, à des styles CSS ou SCSS, à des imports, à des services injectés et parfois à des modèles de données. Le routage associe ensuite des chemins applicatifs à ces composants. Une migration doit donc traiter ces éléments comme un ensemble cohérent, et non comme des fichiers indépendants.

La figure~\ref{fig:angular-architecture-migration} illustre cette organisation et met en évidence deux périmètres candidats de migration : un premier \textit{Remote} regroupant les éléments liés aux composants A et B, et un second centré sur le composant C. Elle distingue également les services spécifiques, les services partagés, les modèles/types communs, les dépendances UI mutualisées et les interactions avec l'API/backend.

\begin{figure}[H]
\centering
\includegraphics[width=\textwidth]{figures/angular_architecture_migration_updated.png}
\caption{Architecture d'une application Angular à analyser avant une migration vers des micro-frontends}
\label{fig:angular-architecture-migration}
\end{figure}

Ainsi, l'analyse d'une application Angular ne doit pas se limiter à la structure des dossiers. Elle doit identifier les dépendances existantes, les services spécifiques et partagés, les modèles communs et les ressources UI mutualisées, afin de définir des périmètres de \textit{Remotes} cohérents et de limiter les erreurs lors de la migration.


\subsection{Large Language Models}

Un \textit{Large Language Model} (LLM) est un modèle de langage de grande taille entraîné sur de vastes corpus textuels afin d'apprendre des régularités linguistiques, sémantiques et contextuelles. Les LLM modernes reposent majoritairement sur l'architecture Transformer, introduite par Vaswani et al., qui utilise des mécanismes d'attention pour modéliser les relations entre les tokens d'une séquence \cite{vaswani_attention_2017}. À grande échelle, ces modèles peuvent réaliser des tâches variées à partir d'instructions ou d'exemples fournis dans le prompt, comme l'ont montré Brown et al. avec GPT-3 et l'apprentissage \textit{few-shot} \cite{brown_language_models_2020}. Dans un contexte d'ingénierie logicielle, nous pouvons les mobiliser pour assister l'analyse, la génération de code, la documentation ou la reformulation, mais leurs sorties doivent rester contrôlées par des règles, des tests et des validations, car elles peuvent contenir des erreurs ou des propositions non conformes au projet.

\subsection{Agentic AI}

Le concept d'\textit{Agentic AI} désigne des systèmes d'intelligence artificielle capables de poursuivre un objectif en combinant compréhension du contexte, raisonnement, planification, utilisation d'outils et production d'actions dans un environnement. Cette notion prolonge celle d'agent intelligent, défini dans la littérature comme une entité située dans un environnement et capable d'agir de manière autonome pour atteindre des objectifs \cite{wooldridge_jennings_intelligent_agents_1995,wooldridge_introduction_2009}. Dans les systèmes récents, un LLM peut jouer le rôle de noyau de raisonnement, tandis que des outils externes permettent d'observer l'état du système, de lire des fichiers, de générer des artefacts ou de vérifier un résultat. Des approches comme ReAct montrent l'intérêt d'alterner raisonnement et actions afin de mieux contrôler les tâches complexes \cite{yao_react_2023}. Dans notre projet, cette logique justifie la séparation entre agents spécialisés, état partagé, orchestrateur et validateurs.

\subsection{Automatisation de la migration}

L'automatisation de la migration vise à exécuter de manière reproductible les tâches répétitives : extraction d'informations, génération de fichiers, adaptation d'imports, réorganisation des routes et vérification des artefacts produits. Elle ne remplace pas la décision d'architecture, mais elle réduit l'effort manuel et rend les transformations plus traçables.

Plusieurs techniques sont mobilisables. L'analyse statique du code, éventuellement fondée sur l'AST, permet d'identifier les composants, services, routes, imports et dépendances sans exécuter l'application \cite{typescript_compiler_api,ts_morph}. La génération par modèles produit des structures répétitives, comme les fichiers de configuration ou de routage \cite{angular_schematics,nx_generators}. Les transformations de code, ou \textit{codemods}, permettent ensuite d'appliquer des modifications à grande échelle, par exemple sur les chemins d'importation ou les déclarations de modules \cite{jscodeshift,openrewrite}.

Les modèles de langage peuvent compléter ces techniques pour proposer un regroupement fonctionnel ou produire une documentation lisible. Ils doivent toutefois être encadrés par des règles et des validateurs, car les décisions critiques doivent rester vérifiables. LangGraph s'inscrit dans cette logique lorsqu'il est utilisé pour orchestrer un workflow avec état, transitions contrôlées et possibilités de reprise \cite{langgraph_docs,langgraph_durable_execution}.

\subsection{Systèmes multi-agents et orchestration}

Un système multi-agent répartit un problème complexe entre plusieurs agents spécialisés. Dans le cas d'une migration automatisée, cette approche permet de distinguer les tâches d'analyse, de conception, de génération, de transformation et de validation, au lieu de les concentrer dans un traitement unique.

Dans son ouvrage \textit{Building Applications with AI Agents: Designing and Implementing Multiagent Systems}, Michael Albada présente notamment le modèle, les outils, la mémoire et l'orchestration comme des éléments structurants d'une application agentique \cite{albada_agents}. Ces notions sont utiles pour concevoir un workflow de migration : les agents ont besoin d'outils pour analyser ou générer, d'une mémoire partagée pour conserver les artefacts intermédiaires et d'un mécanisme d'orchestration pour contrôler l'ordre des étapes.

L'orchestration est particulièrement importante lorsque le processus n'est pas strictement linéaire. Un graphe de workflow permet de représenter les étapes, les transitions conditionnelles, les arrêts en cas d'erreur et les reprises éventuelles. Cette organisation prépare le passage vers le chapitre suivant, où l'architecture proposée détaille concrètement l'orchestrateur, l'état partagé, les agents et les validateurs du pipeline.

\paragraph{Synthèse.}
Les notions présentées dans cette section constituent le socle théorique du projet. Elles expliquent pourquoi la migration d'un monolithe Angular vers des micro-frontends nécessite une analyse structurée, une automatisation contrôlée, une orchestration multi-agents et des mécanismes de validation progressive.


\section{Positionnement par rapport aux solutions existantes}

Dans cette section, nous présentons le positionnement du projet par rapport aux solutions existantes proches du besoin étudié. Cette comparaison permet de montrer les limites des approches disponibles et de préciser la place de la solution proposée dans le contexte de la migration vers une architecture micro-frontend.

Avant de présenter la solution proposée, il est utile de situer le projet par rapport à quelques outils proches du besoin étudié. Les solutions comparées dans le tableau~\ref{tab:etude_existant_microfrontend} couvrent soit l'intégration de micro-frontends, soit l'automatisation de transformations de code, mais elles ne répondent pas entièrement au besoin de migration automatisée d'un projet Angular monolithique vers une architecture \textit{Shell}/\textit{Remotes}.

\begingroup
\small
\renewcommand{\arraystretch}{1.2}
\setlength{\tabcolsep}{3pt}
\begin{longtable}{|>{\raggedright\arraybackslash}p{2.2cm}|>{\raggedright\arraybackslash}p{3.05cm}|>{\raggedright\arraybackslash}p{3.05cm}|>{\raggedright\arraybackslash}p{3.05cm}|>{\centering\arraybackslash}p{1.7cm}|}
\caption{Étude comparative des solutions existantes proches du besoin}\label{tab:etude_existant_microfrontend}\\
\hline
\textbf{Solution} & \textbf{Principe} & \textbf{Apport principal} & \textbf{Limite par rapport au projet} & \textbf{Adéquation} \\
\hline
\textbf{Luigi} \cite{luigi} & Framework destiné à la construction d'applications modulaires et micro-frontends. & Fournit un cadre d'intégration pour composer plusieurs parties applicatives. & Ne prend pas en charge l'analyse complète d'un projet Angular source ni la génération d'un projet migré. & Forte \\
\hline
\textbf{Piral} \cite{piral_migration} & Framework micro-frontend orienté composition, extension et réutilisation de parties applicatives. & Facilite la modularité et la migration progressive dans son propre écosystème. & Plus naturellement adapté à des contextes React et fortement lié à son écosystème. & Moyenne \\
\hline
\textbf{OpenRewrite} \cite{openrewrite} & Framework d'automatisation de refactoring et de transformation de code à grande échelle. & Permet d'industrialiser des transformations répétitives. & Ne couvre pas la logique architecturale complète d'une migration vers un \textit{Shell} et des \textit{Remotes}. & Faible \\
\hline
\end{longtable}
\endgroup

Cette comparaison montre que le besoin du projet se situe entre deux familles d'outils : les frameworks d'intégration micro-frontend et les outils de transformation de code. La solution proposée cherche donc à combiner analyse, génération, migration et validation progressive dans un même processus contrôlé.

\section{Solution proposée}

Dans cette section, nous présentons la solution proposée pour répondre à la problématique du projet. Nous décrivons son principe général, son organisation sous forme de pipeline multi-agent ainsi que les responsabilités principales nécessaires à la migration automatisée.

Pour répondre à la problématique, nous proposons un système multi-agent capable d'assister la migration d'un projet Angular monolithique vers une architecture micro-frontend exécutable. La solution est organisée comme un pipeline : chaque étape reçoit des informations, produit des artefacts intermédiaires et transmet un résultat contrôlé à l'étape suivante.

La solution couvre les principales phases du processus : analyse du projet source, proposition d'un découpage fonctionnel, génération de l'architecture cible, migration des éléments applicatifs essentiels, validation et documentation. Chaque agent prend en charge une responsabilité précise, tandis qu'un mécanisme d'orchestration coordonne l'enchaînement des traitements et les éventuelles corrections.

Le tableau~\ref{tab:responsabilites-solution} présente les responsabilités principales de cette solution et le rôle de chacune dans la chaîne de migration.

\begingroup
\renewcommand{\arraystretch}{1.25}
\begin{longtable}{|p{4cm}|p{10cm}|}
\caption{Vue générale des responsabilités de la solution}\label{tab:responsabilites-solution}\\
\hline
\textbf{Responsabilité} & \textbf{Rôle dans la solution} \\
\hline
Analyse & Comprendre la structure du projet source et produire une représentation exploitable par les autres étapes. \\
\hline
Conception cible & Définir les domaines fonctionnels, les futures applications et les règles d'organisation. \\
\hline
Génération & Créer la structure de base du projet cible et préparer les emplacements nécessaires. \\
\hline
Migration & Transférer les éléments du projet source vers les parties correspondantes du projet cible. \\
\hline
Validation & Contrôler les résultats intermédiaires et signaler les incohérences avant la livraison finale. \\
\hline
Documentation & Expliquer les choix, les fichiers générés et les actions effectuées pendant la migration. \\
\hline
\end{longtable}
\endgroup

L'objectif final est de produire un projet structuré, vérifiable et directement exploitable. La valeur de la solution ne réside donc pas uniquement dans l'automatisation, mais aussi dans la capacité à rendre la migration plus contrôlée, plus explicable et plus simple à reprendre par un développeur.

\section{Méthodologie du travail}

Dans cette section, nous présentons la méthodologie adoptée pour organiser la réalisation du projet. Nous décrivons l'utilisation de Scrum, son application dans le cadre du stage ainsi que les mécanismes de suivi et de validation mis en place.

\subsection{Méthodologie adoptée}

La réalisation du projet s'est appuyée sur Scrum, une méthode agile fondée sur une progression itérative et incrémentale. Le travail est organisé en cycles courts appelés \textit{sprints}, au terme desquels un incrément du produit est conçu, développé et évalué \cite{scrum_guide_fr}.

Ce choix est adapté à la nature du projet. La conception d'un système de migration automatisée comporte plusieurs incertitudes : la qualité de l'analyse du projet source, la pertinence du découpage, la capacité à générer une architecture correcte et la fiabilité des validations. Une démarche itérative permet d'aborder ces difficultés progressivement, en vérifiant régulièrement les résultats obtenus.

\subsection{Application de Scrum au projet}

Dans notre cas, Scrum a été utilisé comme cadre d'organisation pour planifier les tâches, suivre l'avancement et structurer les livrables. Les détails relatifs à l'équipe Scrum, au backlog produit et à la planification des sprints sont présentés dans le chapitre d'analyse et de planification. Ici, nous retenons surtout l'intérêt de cette méthodologie pour organiser un travail évolutif.

La méthodologie a permis de découper le projet en objectifs intermédiaires : préparation de l'application de référence, analyse structurelle, conception cible, génération du projet micro-frontend, migration détaillée, validation et documentation. Chaque objectif a été associé à des tâches réalisables, vérifiables et intégrables dans la suite du pipeline.

La figure~\ref{fig:scrum_process} illustre le déroulement général du processus Scrum et montre la logique itérative adoptée pour organiser le travail.

\begin{figure}[H]
    \centering
    \includegraphics[width=0.75\textwidth]{figures/scrum.png}
    \caption{Déroulement général du processus Scrum \cite{scrum_figure}}
    \label{fig:scrum_process}
\end{figure}

\subsection{Suivi et validation des travaux}

Le suivi du travail a reposé sur la définition de livrables intermédiaires. À chaque étape, le résultat produit devait pouvoir être observé : un analyseur capable d'extraire une structure, un plan de migration lisible, une architecture cible générée, un ensemble d'artefacts migrés ou un rapport de validation. Cette logique a permis d'éviter une progression uniquement basée sur du code non vérifié.

La validation a occupé une place importante dans la démarche. Plutôt que d'attendre la fin du projet pour vérifier l'ensemble de la migration, des contrôles ont été prévus dès les étapes intermédiaires. Cette organisation est cohérente avec l'objectif général du sujet : automatiser la migration sans perdre la maîtrise du résultat.

\section*{Conclusion}
\addcontentsline{toc}{section}{Conclusion}

Ce chapitre a présenté l'entreprise d'accueil, le contexte professionnel du stage, les motivations du sujet, la problématique, les objectifs, le périmètre, les fondements technologiques, le positionnement de la solution et la méthodologie adoptée. Cette organisation permet de regrouper dans un même chapitre le cadre du projet et les notions nécessaires à sa compréhension, tout en évitant les répétitions avec les chapitres d'analyse, de conception et de réalisation.
